% =============================================================================
%   VII. LIMITATIONS AND THREATS TO VALIDITY
% =============================================================================
\section{Limitations and Threats to Validity}
\label{sec:limitations}

\begin{enumerate}
\item \textbf{One window, one cross-section.} The OOS test covers
  \TradingDaysUS trading days in a rising market. With $N=\NOOS$, only
  $|\rho|\ge\MDEOOS$ is detectable at 80\% power. H1 and H3 therefore exclude
  large effects, not small ones, and the risk result (H2) may not generalise
  to other regimes. Extending the window, or repeating the snapshot--hold-out
  cycle, is the most direct next step.
\item \textbf{Hypotheses fixed after the window ended.} The hypotheses and
  analysis code were fixed before OOS prices were downloaded, but after the
  calendar window had passed, and they were not externally registered. Two
  further facts bear on this. First, the profile weights were revised in
  July~2026, inside the window. Second, the authors' general knowledge of
  market conditions could not be excluded. The safeguards are the frozen
  snapshot artefacts and full reporting of alternative weightings (the
  February~2026 weights, and \WShareNeg\% of random weightings, give the
  same return conclusion). The mechanism analyses of
  Section~\ref{sec:res_mechanism} were specified after H1--H3 had been
  evaluated and are exploratory. The second window
  (Section~\ref{sec:prereg_design}) is registered before it opens, so that
  H2 and H4 can be confirmed or refuted there.
\item \textbf{Proxy ESG.} Our E and S pillars are not measurements of
  environmental or social performance (Section~\ref{sec:res_construct}). The
  findings concern public-data ESG as constructed here. They do not
  necessarily transfer to vendor ratings, which have their own validity
  problems \cite{berg_2022, christensen_2022}.
\item \textbf{Governance data dependence.} The governance pillar relies on ISS
  scores redistributed by a free data interface. They are available for
  \NISSFirms of \NMid firms and are sector-imputed otherwise. Their continued
  availability is not guaranteed.
\item \textbf{Universe construction.} The candidate list was assembled from
  index constituents at the time of collection. It is not a point-in-time
  membership history, so earlier-period survivorship cannot be excluded. The
  investability filter and the OOS survivorship bounds address only the
  evaluation window.
\item \textbf{Multi-market pooling.} US and Indian firms are scored against a
  common reference, including \NLarge large caps. Country-neutral returns
  remove market-level differences in outcomes but not differences in
  disclosure regimes. Country subsamples ($N=\NUS$, $N=\NIndia$) are too
  small for separate inference.
\item \textbf{Transaction costs and capacity.} Portfolio results are
  equal-weighted and single-period. We deduct stylised round-trip costs
  (Section~\ref{sec:res_null}), but do not model market impact at scale.
\item \textbf{Provenance split.} The measured/proxy/imputed sub-scores average
  indicator ranks with equal weights, not with the pillar and materiality
  weights of the composite. They show where the signal lives, not how the
  composite would behave with only one data source. The measured cells are
  mostly governance items, so a null for measured cells is not a test of
  measured environmental data. That test requires emissions and workforce
  disclosures (e.g.\ BRSR for the Indian firms), which we leave to future
  work.
\end{enumerate}
